\intermezzo{Overvoltages}{}

\begin{frame}{Overvoltages definition according to IEC}
  \centering
  \topoimg[width=\linewidth,height=0.8\textheight,keepaspectratio]{\icfig{m07_image4}}
\end{frame}

\begin{frame}{Transient overvoltage in the AC grid}
  \begin{columns}[c,onlytextwidth]
    \begin{column}{0.48\textwidth}
      \centering
      \topoimg[width=\linewidth,height=0.8\textheight,keepaspectratio]{\icfig{m07_image5}}
    \end{column}
    \begin{column}{0.48\textwidth}
      \centering
      \topoimg[width=\linewidth,height=0.7\textheight,keepaspectratio]{\icfig{m07_image6}}
    \end{column}
  \end{columns}
  \vspace{0.4em}
  {\centering\color{atGray}\footnotesize IEEE Std 1313.1-1996\par}
\end{frame}

\begin{frame}{Transient overvoltage in the AC grid}
  \begin{columns}[c,onlytextwidth]
    \begin{column}{0.4\textwidth}
      \begin{itemize}
        \setlength{\itemsep}{0.5em}
        \item Overvoltages in the AC grid reflect the ``experience of the
              world'', taking into account protective devices and methods of
              overvoltage limitation.
        \item For AC equipment other overvoltage levels can be specified,
              e.g. transformers.
      \end{itemize}
      \vspace{0.3em}

    \end{column}
    \begin{column}{0.58\textwidth}
      \centering
      \topoimg[width=\linewidth,height=0.66\textheight,keepaspectratio]{\icfig{m07_image7}}
      {\color{atGray}\footnotesize IEEE Std C57.12.00\texttrademark-2021}      
    \end{column}
  \end{columns}
\end{frame}

\begin{frame}{Temporary overvoltage in the AC grid}
  \begin{columns}[c,onlytextwidth]
    \begin{column}{0.54\textwidth}
      \begin{itemize}
        \setlength{\itemsep}{0.5em}
        \item Temporary overvoltages are defined by grid codes such as IEEE 2800, NERC PRC-024 or NERC PRC-029
        \item Temporary overvoltages do not have an impact on the insulation coordination in the AC grid.
      \end{itemize}
      \vspace{0.3em}
    \end{column}
    \begin{column}{0.42\textwidth}
      \centering
      \topoimg[width=\linewidth,height=0.66\textheight,keepaspectratio]{\icfig{m07_image8}}
      {\color{atGray}\footnotesize PRC-024-3 - Frequency and Voltage Protection Settings for Generating Resources.}
    \end{column}
  \end{columns}
\end{frame}

\begin{frame}{Insulation strength}
  \begin{columns}[c,onlytextwidth]
    \begin{column}{0.46\textwidth}
      Insulation strength can be verified by:
      \begin{itemize}
        \setlength{\itemsep}{0.4em}
        \item Testing in a high-voltage lab
        \item Definition of air clearances
      \end{itemize}
      \vspace{0.3em}
      {\color{atGray}\footnotesize www.highvolt.com}
    \end{column}
    \begin{column}{0.50\textwidth}
      \centering
      \topoimg[width=\linewidth,height=0.64\textheight,keepaspectratio]{\icfig{m07_image9}}
    \end{column}
  \end{columns}
\end{frame}

\begin{frame}{Air clearances}
  \begin{columns}[c,onlytextwidth]
    \begin{column}{0.48\textwidth}
      \centering
      \topoimg[width=\linewidth,height=0.66\textheight,keepaspectratio]{\icfig{m07_image11}}
    \end{column}
    \begin{column}{0.48\textwidth}
      \centering
      \topoimg[width=\linewidth,height=0.66\textheight,keepaspectratio]{\icfig{m07_image12}}
    \end{column}
  \end{columns}
  \vspace{0.3em}
  {\color{atGray}\footnotesize IEC 60071.}
\end{frame}

\begin{frame}{Environmental conditions}
  Overvoltages apply for ``normal environmental conditions'' and ``standard
  reference atmospheric conditions'' (IEC 60071-1, §5.9).
  \begin{columns}[T,onlytextwidth]
    \begin{column}{0.48\textwidth}
      \textbf{Normal environmental conditions}
      \begin{itemize}
        \setlength{\itemsep}{0.3em}
        \item Temperature range $-40$\,\textdegree C\,/\,$+40$\,\textdegree C
        \item Altitude $<1000$\,m
        \item Ambient air not significantly polluted
      \end{itemize}
    \end{column}
    \begin{column}{0.48\textwidth}
      \textbf{Standard reference atmosphere}
      \begin{itemize}
        \setlength{\itemsep}{0.3em}
        \item Temperature: 20\,\textdegree C
        \item Pressure: $p_0 = 1013$\,hPa
        \item Absolute humidity: $h_0 = 11$\,g/m$^3$
      \end{itemize}
    \end{column}
  \end{columns}
  \vspace{0.4em}
  {\color{atGray}\footnotesize Outdoors, temperature and absolute humidity
  compensate each other (no correction required); indoors, correction per
  IEC 60060-1 is required.}
\end{frame}

\begin{frame}{Overvoltage definition}
  \begin{columns}[c,onlytextwidth]
    \begin{column}{0.54\textwidth}
      \begin{itemize}
        \setlength{\itemsep}{0.45em}
        \item Overvoltages inside the converter are not based on experience
              values.
        \item Definition must be based on EMT-type simulations.
        \item Component modelling must be adjusted to the anticipated analysis.
        \item Simulation time step must be appropriate to the simulation.
      \end{itemize}
    \end{column}
    \begin{column}{0.42\textwidth}
      \centering
      \topoimg[width=\linewidth,height=0.6\textheight,keepaspectratio]{\icfig{m07_image13}}\\[0.3em]
      {\color{atGray}\footnotesize Transformer representation considering
      capacitive effects.}
    \end{column}
  \end{columns}
\end{frame}

\begin{frame}{Faults for overvoltage definition}
  \begin{columns}[c,onlytextwidth]
    \begin{column}{0.54\textwidth}
      \begin{itemize}
        \setlength{\itemsep}{0.45em}
        \item Overvoltages within the converter are defined by converter
              internal faults.
        \item Converter internal faults are given by the converter design.
        \item Fault clearance depends on: DC protection (DC fault detection)
              and circuit-breaker operation.
      \end{itemize}
      \vspace{0.3em}
      {\color{atGray}\footnotesize CIGRE 832 --- Guide for Electromagnetic
      Transient Studies Involving VSC Converters.}
    \end{column}
    \begin{column}{0.42\textwidth}
      \centering
      \topoimg[width=\linewidth,height=0.66\textheight,keepaspectratio]{\icfig{m07_image14}}
    \end{column}
  \end{columns}
\end{frame}
